\textbf{Proof.}
Write uniform asymptotic coverage as
\[
 \inf_{Q\in\mathcal M_{n,m}(a,s_1,t)\cup
                   \mathcal M_{n,m}(a,s_2,t)}
 Q\{\tau(Q)\in\mathcal C_{n,m}\}
 \ge1-\gamma-o(1).
\]
We first record the interval form of the two-point argument.  If \(Q_0,Q_1\) have parameter separation \(\Delta\), then an interval containing both parameter values has length at least \(\Delta\).  Consequently,
\[
\begin{aligned}
 \mathbb E_{Q_0}|\mathcal C_{n,m}|
 &\ge \Delta\,Q_0\{\tau(Q_0),\tau(Q_1)\in\mathcal C_{n,m}\}\\
 &\ge \Delta\bigl[Q_0\{\tau(Q_0)\in\mathcal C_{n,m}\}
       +Q_1\{\tau(Q_1)\in\mathcal C_{n,m}\}
       -1-\operatorname{TV}(Q_0,Q_1)\bigr]\\
 &\ge \Delta\{1-2\gamma-\operatorname{TV}(Q_0,Q_1)-o(1)\}.
\end{aligned}
\]
Here the second line uses
\(Q_0(B)\ge Q_1(B)-\operatorname{TV}(Q_0,Q_1)\).

We construct both alternatives around one smoother reference law.  Put
\[
 f_S(x)=(a+1)x^a,\qquad f_T(x)=1,
 \qquad \ell(x)=x-\frac12.
\]
Choose a fixed sufficiently small \(\delta_0>0\), and take Bernoulli potential outcomes with
\[
 \mu_0(x)=\frac12,
 \qquad
 \mu_1(x)=\frac12+\delta_0\ell(x).
\]
Use the fixed known propensity, make treatment conditionally independent of the potential outcomes given \(X\), impose consistency, and give the target potential outcomes the same conditional means.  The constant density factors satisfy every order-\(t\) condition.  Because \(\ell\) is smooth, reducing \(\delta_0\) puts both responses strictly inside the radius-\(M\) balls for orders \(s_1\) and \(s_2\), keeps their values in \([0,1]\), and ensures residual variance at least \(M^{-1}\).  Denote this observed law by \(Q_0\).  Then
\(Q_0\in\mathcal M_{n,m}(a,s_2,t)\), and it also satisfies the order-\(s_1\) response restrictions.

Set
\[
 h=n^{-1/(2s_1+a+1)},
 \qquad
 b_{h,s_1}(x)=h^{s_1}K(x/h)\mathbf 1\{0\le x\le h\}.
\]
Construct \(Q_1\) by replacing only the treated response by
\[
 \mu_{1,1}(x)=\frac12+\delta_0\ell(x)+\delta b_{h,s_1}(x).
\]
The boundary-bump scaling and a sufficiently small fixed \(\delta>0\) put
\(Q_1\) in \(\mathcal M_{n,m}(a,s_1,t)\).  The baseline Bernoulli probabilities are uniformly separated from zero and one, so the calculation in
\(\mathrm{lem:boundary\text{-}bump}\) gives
\[
 \operatorname{KL}(Q_1,Q_0)
 \le Cn\delta^2h^{2s_1+a+1}
 =C\delta^2,
\]
whereas
\[
 \Delta_S:=|\tau(Q_1)-\tau(Q_0)|
 =\delta h^{s_1+1}\int_0^1K(u)\,du
 \ge c_Sn^{-(s_1+1)/(2s_1+a+1)}.
\]
By the log-sum calculation
\(\operatorname{TV}(Q_1,Q_0)\le
\sqrt{\operatorname{KL}(Q_1,Q_0)/2}\), reducing \(\delta\) ensures
\(\operatorname{TV}(Q_1,Q_0)\le\eta_S<1-2\gamma\).  The interval inequality then yields, for all sufficiently large indices,
\[
 \mathbb E_{Q_0}|\mathcal C_{n,m}|
 \ge c_1n^{-(s_1+1)/(2s_1+a+1)}.
\]

For the target alternative, retain exactly the same source law and responses and change only the target density to
\[
 f_{T,2}(x)=1+u_m\ell(x),
 \qquad u_m=\delta_1m^{-1/2}.
\]
The equality \(\int_0^1\ell=0\), strict target-factor slack, and smoothness of \(\ell\) make this an admissible density for sufficiently small fixed \(\delta_1>0\).  Let \(Q_2\) be the resulting law, with target potential outcomes chosen to preserve mean transportability.  It belongs to the smoother class.  Since the source laws are identical and the reference target density is one,
\[
 \operatorname{KL}(Q_2,Q_0)
 \le m u_m^2\int_0^1\ell(x)^2\,dx
 \le C\delta_1^2,
\]
while
\[
 \Delta_T:=|\tau(Q_2)-\tau(Q_0)|
 =\delta_0u_m\int_0^1\ell(x)^2\,dx
 =\frac{\delta_0\delta_1}{12}m^{-1/2}.
\]
Reduce \(\delta_1\) so that
\(\operatorname{TV}(Q_2,Q_0)\le\eta_T<1-2\gamma\).  Applying the same interval inequality under the identical reference law \(Q_0\) gives
\[
 \mathbb E_{Q_0}|\mathcal C_{n,m}|
 \ge c_2m^{-1/2}.
\]
Both lower bounds hold at the same law
\(Q_0\in\mathcal M_{n,m}(a,s_2,t)\).  Finally,
\(\max\{u,v\}\ge(u+v)/2\) yields
\[
 \mathbb E_{Q_0}|\mathcal C_{n,m}|
 \ge c\left\{m^{-1/2}
 +n^{-(s_1+1)/(2s_1+a+1)}\right\}.
\]
Thus honest adaptation over the two fixed H\"older classes forces the rough-class boundary rate at a smoother law. \(\square\)